\scanpage{44}
$\varphi_g=0$ i.e. $v(g(P))=0$ pour tt $P\in X(\overline K)$, $v\in M_{\overline K}$, donc $g(P)$ doit être une $M_{\overline K}$-unité pour tout $P\in X(\overline K)$.

Si $g$ n'est pas \unsure{constante|constant} sur les composantes \deleted{connexes} \inserted{irréd.} géom. de $X$, alors \deleted{\unsure{$g$}} on aurait $g(X(\overline K))=\overline K-\text{ens fini}$, donc \deleted{on aurait pour $v\in M_{\overline K}$} comme $v$ \inserted{$\in M_{\overline K}$} prend une infinité de valeurs, la propriété au-dessus ne serait pas vérifiée. Ainsi $g$ est algébrique sur $K$, et de plus en tant qu'\unsure{él\textsuperscript{mt}|élément|dit} de \inserted{la $K$-}algèbre séparable des sections \deleted{séparables} de $\underline O_X^{\times}$ alg. sur $K$, cela doit être une $M_K$-\uline{unité}.

Dans le cas particulier où $X$ est géom. intègre, on voit donc que le noyau de $\underline R_X^{\times}\unsure{-|\to}\Phi_X^*$ est le faisceau \deleted{\unsure{qui}} constant $U_X$ défini par le groupe $U$ des $M_K$-unités de $K$. On a donc les suites exactes can.
% archive sheet number: 43
